% BORRADOR (comentarios del director, 19-oct-2026): reemplaza MainMatter/Cap4/M05-Chapter4.tex. Borrar esta línea al pasarlo.
Este capítulo sintetiza el artefacto evaluado: la arquitectura de JobAI, su pipeline de consulta en la configuración de referencia, que es la que se evalúa, y las dos iteraciones de diseño de este trabajo. Los detalles de implementación están en el Anexo~\ref{anx:artefacto}, y las fallas que la evaluación reveló durante las iteraciones, en la Sección~\ref{sec:res-fallas}.

\section{Arquitectura}
\label{sec:art-arquitectura}

JobAI es una plataforma de agregación de ofertas de empleo remoto cuyo componente central es un copiloto conversacional basado en RAG. Se organiza en tres capas (Figura~\ref{fig:art-arquitectura}): un cliente web en Next.js, que atiende la autenticación, la navegación del catálogo y el panel conversacional; un servidor en NestJS, que orquesta la ingesta, la recuperación y la generación, y una capa de datos con PostgreSQL como almacén relacional y un almacén vectorial que, desde la Iteración~1, puede ser pgvector o Qdrant. Un servidor propio, ejecutado localmente sobre una GPU, sirve el modelo de \textit{embeddings} y el de reordenamiento; los modelos de lenguaje se consumen a través de un proveedor externo, Groq. Los módulos del servidor se describen en el Anexo~\ref{anx:artefacto}.

El capítulo describe la \emph{configuración de referencia} del artefacto, definida en el repositorio junto con el código. El despliegue público de la plataforma conserva todavía la configuración anterior a la Iteración~2, con un proveedor externo de \textit{embeddings} y sin reordenamiento.

\begin{figure}[htbp!]
    \centering
    \caption{Arquitectura de JobAI}
    \label{fig:art-arquitectura}
    \begin{tikzpicture}[
        box/.style={draw, rounded corners, align=center, minimum height=1cm, text width=3.2cm, font=\small},
        store/.style={draw, cylinder, shape border rotate=90, aspect=0.15, align=center, minimum height=1.2cm, text width=2.4cm, font=\small},
        arr/.style={->, >=stealth}
    ]
        \node[box, text width=2.2cm] (fe) at (0,0) {Cliente web\\(Next.js)};
        \node[box, text width=2.2cm] (be) at (5.4,0) {Servidor\\(NestJS)};
        \node[box] (models) at (10.2,2.1) {Modelos de lenguaje\\(Groq)};
        \node[box] (local) at (10.2,0) {\textit{Embeddings} y reordenamiento (servidor local, GPU)};
        \node[box] (sources) at (10.2,-2.1) {Fuentes de ofertas\\(Remotive, Jobicy, Findwork)};
        \node[store] (pg) at (3.8,-3.2) {PostgreSQL\\(+ pgvector)};
        \node[store] (qd) at (7.0,-3.2) {Qdrant};
        \draw[arr] (fe) -- node[above, font=\scriptsize] {REST / NDJSON} (be);
        \draw[arr] (be) -- (models);
        \draw[arr] (be) -- (local);
        \draw[arr] (sources) -- (be);
        \draw[arr] (be) -- (pg);
        \draw[arr, dashed] (be) -- (qd);
    \end{tikzpicture}
\end{figure}

\section{Pipeline de consulta}
\label{sec:art-consulta}

La Figura~\ref{fig:art-pipeline} resume el recorrido de una consulta y la Tabla~\ref{tab:art-config}, los parámetros de la configuración de referencia.

\begin{figure}[htbp!]
    \centering
    \caption{Pipeline de consulta de JobAI}
    \label{fig:art-pipeline}
    \begin{tikzpicture}[
        stage/.style={draw, rounded corners, align=center, text width=4.2cm, minimum height=0.9cm, font=\small},
        half/.style={stage, text width=3.2cm},
        decision/.style={draw, diamond, aspect=2.2, align=center, inner sep=1pt, font=\small},
        arr/.style={->, >=stealth}
    ]
        \node[stage] (q) at (0,0) {Consulta del usuario};
        \node[decision] (cls) at (0,-1.6) {Clasificación};
        \node[stage] (rag) at (-3.6,-3.2) {Ruta de recuperación\\(con restricciones como filtro)};
        \node[stage] (agg) at (3.6,-3.2) {Ruta de agregación\\(plantilla SQL)};
        \node[half] (dense) at (-5.5,-4.8) {Búsqueda densa: 50 fragmentos $\rightarrow$ ofertas};
        \node[half] (bm25) at (-1.7,-4.8) {BM25:\\50 ofertas};
        \node[stage] (rrf) at (-3.6,-6.4) {Fusión RRF\\(20 candidatas)};
        \node[stage] (rr) at (-3.6,-7.8) {Reordenamiento\\(5 ofertas)};
        \node[stage] (ctx) at (-3.6,-9.2) {Contexto por oferta + campos estructurados};
        \node[stage] (sum) at (3.6,-4.8) {Resumen de la agregación};
        \node[stage] (llm) at (0,-10.8) {Prompt + generación (LLM)};
        \node[stage] (out) at (0,-12.2) {Respuesta incremental con fuentes};
        \draw[arr] (q) -- (cls);
        \draw[arr] (cls) -| node[above, pos=0.25, font=\scriptsize] {recuperación} (rag);
        \draw[arr] (cls) -| node[above, pos=0.25, font=\scriptsize] {agregación} (agg);
        \draw[arr] (rag) -- (dense);
        \draw[arr] (rag) -- (bm25);
        \draw[arr] (dense) -- (rrf);
        \draw[arr] (bm25) -- (rrf);
        \draw[arr] (rrf) -- (rr);
        \draw[arr] (rr) -- (ctx);
        \draw[arr] (agg) -- (sum);
        \draw[arr] (ctx) |- (llm);
        \draw[arr] (sum) |- (llm);
        \draw[arr] (llm) -- (out);
    \end{tikzpicture}
\end{figure}

\begin{table}[htbp!]
    \centering
    \caption{Configuración de referencia del artefacto}
    \label{tab:art-config}
    \small
    \begin{tabular}{|p{4.4cm}|p{9.8cm}|}
        \hline
        \textbf{Componente} & \textbf{Configuración} \\
        \hline
        Modelo de \textit{embeddings} & multilingual-e5-base \citep{Wang2024mE5}, 768 dimensiones, ejecutado localmente; prefijos \texttt{query:} y \texttt{passage:}. e5-base-v2 \citep{Wang2022E5} se conserva para medir el efecto del cambio de modelo \\
        \hline
        Fragmentación & Estructurada (identidad, requisitos, responsabilidades, beneficios y salario); ventana deslizante de 800 caracteres con solapamiento de 200 cuando el análisis estructurado falla \\
        \hline
        Restricciones & Extraídas por el clasificador (salario, modalidad remota, tipo de contrato) y aplicadas como filtro en la recuperación \\
        \hline
        Primera etapa & 50 fragmentos más similares por coseno, agrupados por oferta, y 50 ofertas por BM25, fusionadas con RRF ($k = 60$) \\
        \hline
        Reordenamiento & Qwen3-Reranker-0.6B \citep{Zhang2025Qwen3} sobre las 20 primeras candidatas; se conservan 5 ofertas \\
        \hline
        Índice vectorial & HNSW con $m=16$ y $ef_{construction}=128$ en pgvector y en Qdrant; $ef_{search}=40$ \\
        \hline
        Modelo generador & \texttt{qwen/qwen3.8-27b} servido por Groq, temperatura 0,2 (0 durante la evaluación) \\
        \hline
        Clasificación de consultas & Patrones \textit{regex} más un clasificador basado en LLM restringido a un catálogo cerrado de plantillas SQL \\
        \hline
        Equipo de evaluación & AMD Ryzen 5 3600 (6 núcleos, 12 hilos), 16 GB de RAM y GPU NVIDIA GeForce RTX 3060 de 12 GB, que ejecuta el modelo de \textit{embeddings} y el de reordenamiento \\
        \hline
    \end{tabular}
\end{table}

\paragraph{Ingesta.}
\label{sec:art-ingesta}
Un proceso programado consulta las tres fuentes públicas cada seis horas y puede desactivarse, como se hizo para congelar el corpus de evaluación. Deduplica las ofertas, extrae de cada una con un modelo de lenguaje sus campos estructurados (resumen, salario, responsabilidades, requisitos, beneficios, habilidades, tipo de contrato y modalidad), normaliza los salarios a dólares anuales y la divide en fragmentos \emph{tipados}: identidad, requisitos, responsabilidades y habilidades, y beneficios y salario. La fragmentación tipada responde a un hallazgo de una iteración previa (Sección~\ref{sec:met-ciclos-previos}): un fragmento que contiene solo la información por la que se pregunta compite mejor en la recuperación que uno que la mezcla con otra. La configuración de referencia usa \texttt{gpt-oss-120b} para el análisis estructurado, pero el corpus de evaluación se procesó con Llama 3.1 8B ejecutado localmente (Sección~\ref{sec:met-corpus}).

\paragraph{Clasificación de la consulta.}
\label{sec:art-clasificacion}
Patrones \textit{regex} detectan señales de recuperación, de agregación y de restricción, en inglés y en español. Una consulta con señales de recuperación y sin las otras va directamente a la recuperación, sin invocar ningún modelo de lenguaje. En los demás casos, un clasificador basado en LLM, a temperatura 0, elige entre recuperación, agregación o una ruta híbrida y, según el caso, selecciona una de 24 plantillas SQL parametrizadas o extrae las restricciones de la consulta: salario mínimo o máximo, modalidad remota y tipo de contrato. El modelo nunca genera SQL libre. Tres reglas deterministas corrigen la salida del clasificador cuando su etiqueta y su plantilla se contradicen (Anexo~\ref{anx:clasificacion}). El título del cargo y las ubicaciones distintas de \textit{remoto} no se convierten en filtros: exigir su coincidencia literal reproduciría la limitación léxica que el pipeline busca superar.

\paragraph{Recuperación.}
\label{sec:art-recuperacion}
La búsqueda densa recupera los 50 fragmentos más similares que cumplen el filtro y los agrupa por oferta, con la similitud del mejor fragmento; si el usuario cargó un currículo, antepone a la consulta un perfil breve. En paralelo, BM25 \citep{Robertson2009} busca las 50 ofertas que mejor coinciden con el texto de la consulta bajo el mismo filtro. Las dos listas se combinan mediante fusión por rangos recíprocos (RRF) \citep{Cormack2009}, que suma $1/(k + r)$ por cada lista en la que la oferta aparece en la posición $r$, con $k = 60$, sin necesidad de calibrar entre sí los puntajes. Las 20 primeras ofertas pasan por un modelo de reordenamiento de tipo \textit{cross-encoder}, Qwen3-Reranker-0.6B \citep{Zhang2025Qwen3}, que lee juntas la consulta y la oferta \citep{Nogueira2019} y estima la probabilidad de que la oferta sea relevante; se conservan las cinco mejores. El filtro se aplica dentro de cada motor, durante la búsqueda; si ninguna oferta lo cumple, la búsqueda se repite sin él y el prompt pide advertirlo al usuario. Una compuerta de \emph{sin resultados} responde que no hay ofertas relevantes si la mejor oferta reordenada tiene una probabilidad menor que 0,134, un umbral calibrado con consultas de desarrollo (Sección~\ref{sec:art-iter-recuperacion}). Si el servidor de reordenamiento no responde, se conserva el orden de la fusión.

\paragraph{Contexto y generación.}
Para cada una de las cinco ofertas se arma un bloque de contexto con su encabezado y el texto de sus fragmentos. Si la consulta pregunta por un campo, como el salario o los requisitos, que no aparece en los fragmentos recuperados, se añade desde los datos estructurados de la oferta, para que el modelo no lo complete con información plausible pero inexistente. El prompt instruye al modelo a responder \emph{únicamente} a partir de las ofertas, a indicar cuándo no contienen la información solicitada, a no inventar detalles y a responder en el idioma de la pregunta. El historial reciente de la conversación, hasta 12 mensajes, se incluye solo en la generación; la recuperación usa la consulta actual. En la ruta híbrida, la recuperación y la agregación se ejecutan en paralelo y se combinan en un único prompt.

\section{Iteración 1: corrección}
\label{sec:art-vectorial}

La primera iteración partió de una revisión del artefacto frente a su documentación, hecha al preparar la evaluación, para que lo medido correspondiera a lo que el artefacto dice hacer. Corrigió las fallas que clasifica la Sección~\ref{sec:res-fallas}, la mayoría en la ingesta, la agregación y la búsqueda vectorial, y rediseñó el almacenamiento vectorial para que pgvector y Qdrant pudieran compararse en igualdad de condiciones.

\paragraph{Diagnóstico.} En su versión inicial, JobAI almacenaba los vectores en una columna de PostgreSQL mediante la extensión pgvector \citep{Pgvector}, con un índice IVFFlat creado en la migración inicial, sobre una tabla todavía vacía. La consulta de similitud ordenaba los resultados por la expresión derivada $1 - d(\mathbf{q}, \mathbf{v})$, y el planificador de PostgreSQL solo puede servir una búsqueda de vecinos cercanos desde el índice cuando la consulta ordena directamente por el operador de distancia. La inspección con \texttt{EXPLAIN ANALYZE} sobre el corpus de evaluación (Tabla~\ref{tab:art-explain}) confirmó que la consulta original nunca usaba el índice. También mostró que, con el índice HNSW y la consulta corregida, el planificador sigue prefiriendo el barrido secuencial a esta escala, aunque el recorrido del índice es 12 veces más rápido; por eso la búsqueda admite un modo que fuerza el índice.

\begin{table}[htbp!]
    \centering
    \caption{Plan de ejecución de la búsqueda de 15 vecinos en pgvector sobre el corpus de evaluación}
    \label{tab:art-explain}
    \small
    \begin{tabular}{|p{5.2cm}|p{2.2cm}|p{4.6cm}|r|}
        \hline
        \textbf{Consulta} & \textbf{Índice} & \textbf{Plan elegido} & \textbf{Tiempo} \\
        \hline
        Original (orden por $1-d$) & IVFFlat & Barrido secuencial + ordenamiento & 25,7 ms \\
        \hline
        Reescrita (orden por $d$) & IVFFlat & Recorrido del índice & 2,8 ms \\
        \hline
        Reescrita (orden por $d$) & HNSW & Barrido secuencial + ordenamiento & 24,8 ms \\
        \hline
        Reescrita, índice forzado & HNSW & Recorrido del índice & 2,0 ms \\
        \hline
    \end{tabular}
\end{table}

Este diagnóstico cambia la motivación de la migración. La lentitud observada en pgvector no se debía al motor, sino a una consulta que no usaba el índice. Por eso la migración a Qdrant no se planteó como una mejora supuesta de antemano, sino como una alternativa que se evalúa (PE4).

\paragraph{Diseño de la migración.} Los módulos de ingesta, RAG y gestión de ofertas pasaron a depender de una interfaz única de almacén vectorial, con dos implementaciones seleccionables por configuración. En pgvector, la consulta se reescribió para ordenar por el operador de distancia, el índice IVFFlat se reemplazó por HNSW \citep{Malkov2020} y las búsquedas con filtros activan los barridos iterativos del índice. En Qdrant \citep{QdrantDocs}, la colección usa los mismos parámetros HNSW y guarda como carga útil los atributos que filtran la búsqueda. PostgreSQL sigue siendo la fuente de verdad del texto y de la relación con las ofertas, y un comando idempotente copia los vectores a Qdrant y verifica los conteos (detalles en el Anexo~\ref{anx:migracion}). Sobre el corpus de evaluación, la búsqueda exacta en ambos motores coincidió en el 99,85\,\% de los 10 primeros resultados de 200 consultas con los vectores de e5-base-v2, y en el 99,75\,\% tras recodificar el corpus con el modelo multilingüe; las diferencias son empates entre fragmentos de texto idéntico. Los dos almacenes contienen, por tanto, los mismos datos y usan la misma métrica.

\section{Iteración 2: mejora}
\label{sec:art-iter-recuperacion}

Una primera corrida completa de la evaluación, con la configuración de la Iteración~1, y una revisión del pipeline frente a la literatura mostraron cuatro brechas:
\begin{itemize}
    \item \textbf{Las restricciones de la consulta no llegaban a la búsqueda vectorial} (falla A4). Cuando el clasificador sí las extraía, la plataforma las resolvía con un filtro SQL por subcadena del título, la misma búsqueda léxica que el trabajo pretende superar, y el \textit{benchmark} de búsqueda filtrada (PE4) medía una capacidad que la aplicación no usaba.
    \item \textbf{El pipeline era un RAG \emph{ingenuo}} en la clasificación de \citet{Gao2023}: una única búsqueda densa, sin reordenamiento ni recuperación híbrida.
    \item \textbf{Las consultas en español recibían respuestas en inglés} (falla A6).
    \item \textbf{El modelo de \textit{embeddings} era solo para inglés:} e5-base-v2 \citep{Wang2022E5}, de modo que las consultas en español dependían de la coincidencia accidental de términos.
\end{itemize}

La Tabla~\ref{tab:art-cambios} resume los cambios de la iteración. Los filtros se aplican durante la búsqueda y no sobre los $k$ vecinos ya recuperados, lo que evita listas incompletas con filtros selectivos \citep{Gollapudi2023, Patel2024}. La búsqueda híbrida aporta la coincidencia exacta de términos poco frecuentes, como nombres de tecnologías; se eligió RRF porque no requiere calibrar puntajes ni datos de entrenamiento, con el valor habitual $k = 60$ sin ajustar, aunque \citet{Bruch2023} encontraron que RRF es sensible a ese parámetro y que una combinación convexa de los puntajes la supera. El reordenamiento se limita a 20 candidatas porque su costo crece con su número. multilingual-e5-base conserva la dimensión de e5-base-v2, de modo que el corpus se volvió a codificar en el lugar, y los vectores anteriores se conservaron para medir el efecto del cambio.

\begin{table}[htbp!]
    \centering
    \caption{Cambios de la Iteración~2}
    \label{tab:art-cambios}
    \small
    \begin{tabular}{|p{3cm}|p{5.4cm}|p{5.4cm}|}
        \hline
        \textbf{Aspecto} & \textbf{Antes} & \textbf{Después} \\
        \hline
        Restricciones de la consulta & Ignoradas si la consulta tenía señales de recuperación; si no, filtro SQL por subcadena del título, con 30 ofertas ordenadas por fecha & Extraídas por el clasificador (inglés y español) y aplicadas como filtro dentro de la búsqueda densa y de BM25 \\
        \hline
        Primera etapa & Búsqueda densa de 15 fragmentos con umbral de similitud 0,5 & Búsqueda densa de 50 fragmentos y BM25 de 50 ofertas, fusionadas con RRF \\
        \hline
        Orden final & Similitud del mejor fragmento & Reordenamiento con Qwen3-Reranker-0.6B sobre 20 candidatas \\
        \hline
        Consultas sin ofertas relevantes & Siempre se devolvían cinco ofertas & Compuerta sobre el mejor puntaje del reordenamiento (0,134) \\
        \hline
        Modelo de \textit{embeddings} & e5-base-v2 (inglés) & multilingual-e5-base \\
        \hline
        Idioma de la respuesta & El del contexto & El de la pregunta \\
        \hline
        Salarios fuera de rango & Se almacenaban tal como llegaban & Se registran como no informados \\
        \hline
    \end{tabular}
\end{table}

\subsection{Selección del modelo de reordenamiento}

\begin{table}[!htb]
    \centering
    \caption{Selección del modelo de reordenamiento en el conjunto de desarrollo (28 consultas, 30 candidatas)}
    \label{tab:art-rerank-dev}
    \small
    \begin{tabular}{|l|c|c|}
        \hline
        \textbf{Variante} & \textbf{nDCG@10} & \textbf{P@5} \\
        \hline
        Sin reordenamiento (densa + BM25 + RRF) & 0{,}742 & 0{,}757 \\
        \hline
        bge-reranker-v2-m3 & 0{,}804 & 0{,}814 \\
        \hline
        bge-reranker-v2-m3, fusionado con RRF & 0{,}798 & 0{,}800 \\
        \hline
        Qwen3-Reranker-0.6B & 0{,}820 & 0{,}864 \\
        \hline
        Qwen3-Reranker-0.6B con instrucción & 0{,}839 & 0{,}879 \\
        \hline
        Qwen3-Reranker-0.6B con instrucción, fusionado con RRF & 0{,}812 & 0{,}879 \\
        \hline
    \end{tabular}
\end{table}

El modelo de reordenamiento, el número de candidatas y el umbral de la compuerta se eligieron con el conjunto de desarrollo (Sección~\ref{sec:met-consultas}) y una regla de decisión fijada antes de ver los resultados, para no ajustar el artefacto a las consultas de prueba. En términos de DSR es una evaluación formativa: sus resultados orientan el diseño y no se reportan como resultados del artefacto. Se compararon bge-reranker-v2-m3, basado en BGE-M3 \citep{Chen2024M3}, y Qwen3-Reranker-0.6B, en cinco variantes: con y sin instrucción de tarea, y con el orden dado solo por el reordenamiento o fusionado por RRF con el de la primera etapa. La regla tenía tres etapas:
\begin{enumerate}
    \item Con 30 candidatas, se elige la variante de mayor nDCG@10; si otra queda a 0,02 o menos, se prefiere la más simple.
    \item Se usan 20 candidatas en lugar de 30 si el nDCG@10 baja 0,01 o menos.
    \item El umbral es el mayor valor que deja por debajo, como máximo, a una de las 28 consultas de desarrollo, evaluado sobre el mejor puntaje de cada consulta; se adopta solo si deja por debajo al menos a 8 de las 10 consultas fuera de dominio.
\end{enumerate}


Todas las variantes mejoraron entre 0,06 y 0,10 el nDCG@10 de la recuperación sin reordenamiento (Tabla~\ref{tab:art-rerank-dev}). Qwen3-Reranker-0.6B sin instrucción quedó a 0,019 de la mejor variante, y la regla lo eligió por ser más simple. Pasar de 30 a 20 candidatas cambió el nDCG@10 de 0,832 a 0,830 y redujo la mediana de la latencia de recuperación de 1.184 a 788~ms, así que se adoptaron 20. El mejor puntaje más bajo entre las consultas de desarrollo fue 0,132 y el siguiente 0,134; las diez consultas fuera de dominio quedaron por debajo de 0,09, de modo que el umbral de 0,134 cumplió la condición y se adoptó.
